\section{Parameters}\label{app:params}

\begin{table}[H]
  \centering\small
  \caption{Simulation, sensor, filter and evaluation parameters (defaults; experiments that
    change one say so). Generated from \texttt{configs/default.yaml} and
    \texttt{configs/experiments/\_common.yaml}.}\label{tab:params}
  % Generated by scripts — do not edit.
\begin{tabular}{lllr}
\toprule
Group & Parameter & Value & Unit \\
\midrule
Simulation & Time step $\Delta t$ & 1 & s \\
Simulation & Flight duration & 300 & s \\
Simulation & Ground speed & 40 & m/s \\
Simulation & Heading weave amplitude & 25 & deg \\
Simulation & Heading weave period & 120 & s \\
Simulation & Clearance above highest terrain in corridor & 300 & m \\
Simulation & Corridor half-width & 1000 & m \\
Simulation & Flight margin from map edge / 3 & 2000 & m \\
Sensors & INS initial bias std (per axis) & 0.3 & m/s \\
Sensors & INS bias random walk & 0.002 & m/s/$\sqrt{\mathrm{s}}$ \\
Sensors & INS white velocity noise & 0.05 & m/s \\
Sensors & Barometric altimeter noise & 2 & m \\
Sensors & Radar altimeter noise & 3 & m \\
Sensors & Radar outlier std (E4) & 100 & m \\
Prior & Initial position std $\sigma_0$ (default) & 500 & m \\
Particle filter & Particles $N$ (default) & 5000 &  \\
Particle filter & Resampling scheme & systematic &  \\
Particle filter & Resample when $N_\mathrm{eff} < \tau N$, $\tau$ & 0.5 &  \\
Particle filter & Process noise, position $\sigma_p$ & 2 & m/step \\
Particle filter & Process noise, bias $\sigma_b$ & 0.002 & m/s/step \\
Particle filter & Measurement std $\sigma_z$ & 3.6 & m \\
Particle filter & Roughening $K$ (default) & 0 &  \\
Evaluation & Convergence threshold & 50 & m \\
Evaluation & Tracking window & 100 & s \\
\bottomrule
\end{tabular}

\end{table}

\section{Reproducibility}\label{app:repro}

\paragraph{Machine.} All runtimes refer to one laptop: \MachineCpu{} (6 performance and 4
efficiency cores), 16\,GB RAM, \MachineOs{}, Python \MachinePython{}, NumPy \MachineNumpy{}.

\paragraph{Reproducing everything.} From a fresh clone with an OpenTopography API key in
\texttt{.env}:
\begin{lstlisting}[language=bash,numbers=none]
make setup        # virtual environment + package + test tools
make test         # unit tests
make data         # download, reproject and crop the three regions (4 API calls)
make benchmark    # runtime per filter step (run on an idle machine)
make experiments  # all Monte Carlo experiments (resumable; about an hour on this laptop)
make figures      # every figure, table and number in this report, from saved results
make report       # this PDF (or upload report/ to Overleaf)
\end{lstlisting}
\texttt{make quick} runs every experiment with three runs and short flights as a smoke test, and
\texttt{make smoke} checks the whole pipeline on synthetic terrain without any download.
Each results folder stores the git commit, the data-manifest checksum, the resolved configuration
of every experiment cell and the installed package versions; \texttt{data/manifest.json} stores
SHA-256 checksums of every raster.

\paragraph{Code.} Repository: \texttt{https://github.com/KaveeshwaraBandara/terrain-aided-navigation-pf}.

\paragraph{Data attribution.} Produced using Copernicus WorldDEM-30 \textcopyright{} DLR e.V.
2010--2014 and \textcopyright{} Airbus Defence and Space GmbH 2014--2018 provided under
COPERNICUS by the European Union and ESA; all rights reserved. Data accessed through
OpenTopography.
